\documentclass[12pt,a4paper]{article}
\usepackage[T1]{fontenc}
\usepackage[utf8]{inputenc}
\usepackage[ngerman]{babel}
\usepackage{geometry}
\usepackage{enumitem}
\geometry{margin=2.5cm}
\setlist{nosep}

\title{Vorläufige Gliederung der Bachelorarbeit}
\author{}
\date{}

\begin{document}
\maketitle

\section*{Arbeitstitel}

\textbf{Robustheit von TabPFN-3.5 gegenüber irrelevanten Merkmalen:
Eine kontrollierte Untersuchung im Vergleich zu LightGBM}

\section*{Forschungsfragen}

\textbf{Übergeordnete Forschungsfrage:}
Wie robust ist TabPFN-3.5 gegenüber schrittweise hinzugefügten unabhängigen,
irrelevanten Merkmalen im Vergleich zu LightGBM, und wie ist dieses Ergebnis
anhand der methodischen Entwicklung und bisherigen Evaluationen von TabPFN
einzuordnen?

\medskip
\noindent\textbf{FF1 -- Literaturteil:}
Welche relevanten methodischen Änderungen und adressierten Herausforderungen
tabellarischer Daten beschreiben ausgewählte TabPFN-Arbeiten, und welche
Schlussfolgerungen erlauben ihre Evaluationen?

\medskip
\noindent\textbf{FF2 -- empirischer Teil:}
Wie verändert sich die Klassifikationsleistung von TabPFN-3.5 und LightGBM
bei schrittweiser Zugabe unabhängiger, irrelevanter Merkmale?

Die Versionsentwicklung dient nur der Einordnung von TabPFN-3.5. FF2 bildet
den empirischen Schwerpunkt. Aussagen aus dem Vergleich mit LightGBM werden
nicht auf alle baumbasierten Verfahren übertragen.

\section*{Vorläufige Gliederung}

\begin{enumerate}
  \item \textbf{Einleitung}
    \begin{enumerate}
      \item Problemstellung und Motivation
      \item Zielsetzung und Forschungsfragen
      \item Abgrenzung und Aufbau der Arbeit
    \end{enumerate}
  \item \textbf{Grundlagen}
    \begin{enumerate}
      \item Tabellarische Klassifikationsaufgaben und ihre Besonderheiten
      \item TabPFN als Tabular Foundation Model
      \item Bayesianische Priors und Prior-Data Fitted Networks
      \item Baumbasiertes Gradient Boosting und LightGBM
      \item Irrelevante Merkmale und Robustheit von Klassifikationsmodellen
    \end{enumerate}
  \item \textbf{Stand der Forschung und Literaturanalyse}
    \begin{enumerate}
      \item TabPFN: methodische Entwicklung bis Version 3.5
      \item Auswahl und Einordnung relevanter TabPFN-Evaluationen
      \item Vergleichende Forschung zu Deep Learning und Baumverfahren auf
            Tabellendaten
      \item Bisherige Befunde zu irrelevanten Merkmalen
      \item Forschungsbedarf und Herleitung des Experiments
    \end{enumerate}
  \item \textbf{Methodik}
    \begin{enumerate}
      \item Forschungsdesign und Versuchsbedingungen
      \item Datensätze und Auswahlkriterien
      \item Erzeugung und stufenweise Zugabe irrelevanter Merkmale
      \item Modelle, Konfigurationen und Vorverarbeitung
      \item Datenaufteilungen und Wiederholungen
      \item Metriken und Auswertungsverfahren
    \end{enumerate}
  \item \textbf{Ergebnisse}
    \begin{enumerate}
      \item Datensätze und technische Umsetzung
      \item Ergebnisse ohne zusätzliche irrelevante Merkmale
      \item Veränderung der Leistung bei zunehmender Zahl irrelevanter
            Merkmale
      \item Vergleich der Ergebnisse zwischen Modellen und Datensätzen
    \end{enumerate}
  \item \textbf{Diskussion}
    \begin{enumerate}
      \item Beantwortung der Forschungsfragen
      \item Einordnung in die Literatur
      \item Aussagekraft, Grenzen und Übertragbarkeit
      \item Praktische Bedeutung und weiterer Forschungsbedarf
    \end{enumerate}
  \item \textbf{Fazit}
\end{enumerate}

\section*{Roter Faden}

\begin{quote}
Tabellarische Daten stellen besondere Anforderungen an Lernverfahren
$\longrightarrow$ TabPFN adressiert diese Aufgaben durch einen auf
synthetischen Datensätzen trainierten Prior-Data-Fitted-Ansatz
$\longrightarrow$ bestehende Arbeiten untersuchen die Leistungsfähigkeit
und teilweise die Empfindlichkeit gegenüber irrelevanten Merkmalen
$\longrightarrow$ TabPFN-3.5 erweitert laut technischem Bericht den Umgang
mit größeren Feature-Zahlen, was jedoch nicht automatisch Robustheit
gegenüber gezielt hinzugefügten Störmerkmalen belegt
$\longrightarrow$ die Arbeit untersucht diesen Effekt kontrolliert für
TabPFN-3.5 und LightGBM.
\end{quote}

\end{document}
